\newpage
\chapter*{Introduction Générale}
\addcontentsline{toc}{chapter}{Introduction Générale}
\label{chap:intro-generale}

Le numérique évolue à une vitesse qui ne laisse plus vraiment le temps de s'installer confortablement. Ce qui était considéré comme un avantage concurrentiel hier -- automatiser ses processus, digitaliser ses outils -- est devenu aujourd'hui une simple condition de base. La vraie valeur ajoutée se joue désormais ailleurs : dans la capacité des systèmes à être intelligents, à s'adapter et à prendre en charge des décisions de manière autonome.

Les ressources humaines sont un domaine où ce besoin se fait particulièrement sentir. Recruter, suivre des candidatures, valoriser les profils internes\ldots Ces tâches, bien que fondamentales, restent souvent marquées par des processus lourds et peu cohérents. Des CV non standardisés, des informations incomplètes, des outils qui ne communiquent pas entre eux -- autant de frictions qui font perdre un temps précieux aux équipes RH, au détriment de ce qui compte vraiment : les talents.

Un exemple simple illustre bien ce problème : deux candidats avec des profils similaires peuvent présenter des CV totalement différents en termes de forme. Résultat, la comparaison devient difficile, le traitement s'allonge, et les risques d'erreur d'appréciation augmentent.

C'est ce constat, observé au sein de la société Tritux, qui est à l'origine de ce projet. L'objectif n'était pas simplement d'améliorer l'existant, mais de repenser en profondeur la manière dont les données RH sont collectées, structurées et exploitées.

La solution développée dans ce cadre couvre l'ensemble du cycle RH : centralisation et structuration des profils internes et externes, extraction automatique des informations à partir des CV, standardisation selon des templates aux couleurs de Tritux, assistance intelligente via chatbot, gestion des offres d'emploi synchronisées avec un site public, scoring automatique des candidatures, et automatisation des échanges avec les candidats.

Ce rapport présente les différentes étapes de conception et de réalisation de cette plateforme. Le premier chapitre est consacré à l'étude préalable : contexte, analyse de l'existant, besoins identifiés et méthodologie adoptée. Les chapitres suivants détaillent la mise en œuvre technique, sprint après sprint.

\vfill
